\chapter{Addiction}
\index{Addiction}

\section*{Definition and Overview}
Addiction, also called substance use disorder, is a treatable medical condition in which a person loses control over the use of a substance, such as alcohol, tobacco, prescription medicines, or illicit drugs, despite harmful consequences. It is a chronic brain condition, not a moral failing or a weakness of willpower.

People with addiction may experience intense cravings, continue use even when it causes problems at work, school, or home, and find it very difficult to stop or cut down without help. Addiction is treatable, and many people achieve lasting recovery with the right combination of medical care, counselling, and support.

\section*{Epidemiology}
Addiction is common. In many countries, roughly 1 in 20 adults is dependent on alcohol, and tobacco dependence affects about 1 in 10 adults who smoke daily. Opioids, cannabis, and stimulants account for a large share of drug-related harm, and alcohol and illicit drugs contribute substantially to injury, illness, and premature death.

Addiction can begin at any age but often starts in adolescence or young adulthood. Men are affected more often than women for most substances, although women may progress from first use to dependence more quickly. Rates are higher among people with mental illness, a history of trauma, and social disadvantage.

\section*{Aetiology and Pathophysiology}
Addiction arises from a combination of biological, psychological, and social factors. Substances activate the brain's reward system, releasing dopamine and producing feelings of pleasure or relief. With repeated use, the brain adapts, so that more of the substance is needed to achieve the same effect (tolerance), and stopping causes unpleasant withdrawal symptoms. Over time, changes in brain circuits weaken self-control and strengthen cravings, making drug-seeking behaviour increasingly automatic.

\begin{itemize}
    \item \textbf{Genetic factors:} addiction runs in families
    \item \textbf{Brain chemistry:} altered reward, stress, and impulse-control circuits
    \item \textbf{Psychological factors:} trauma, stress, mental illness, and low mood
    \item \textbf{Social factors:} poverty, peer pressure, and the availability of drugs
    \item \textbf{Early exposure:} starting substance use in adolescence
\end{itemize}

\section*{Clinical Features}
\begin{itemize}
    \item \textbf{Cravings:} strong urges or desires to use the substance
    \item \textbf{Impaired control:} using more, or for longer, than intended
    \item \textbf{Unsuccessful attempts:} difficulty cutting down or stopping despite trying
    \item \textbf{Preoccupation:} spending a great deal of time obtaining, using, or recovering from the substance
    \item \textbf{Neglect:} giving up work, study, family, or hobbies because of use
    \item \textbf{Harm despite consequences:} continuing use even when it causes physical or mental harm
    \item \textbf{Tolerance:} needing more of the substance to get the same effect
    \item \textbf{Withdrawal:} physical or psychological symptoms when the substance is not available
\end{itemize}

\section*{Differential Diagnosis}
Some conditions can look like addiction or accompany it:

\begin{itemize}
    \item Depression, anxiety, and other mental health conditions
    \item Bipolar disorder or other mood disorders
    \item Chronic pain conditions leading to opioid use
    \item Physical (but not harmful) dependence on medicines, which differs from addiction
    \item Encephalopathy or other brain conditions causing poor judgement
\end{itemize}

\section*{When to Seek Medical Care}
Anyone who feels they have lost control over their use of a substance, or whose family is concerned about their use, should speak to a GP. It is especially important to seek help when use is causing harm to health, relationships, finances, or work, or when attempts to stop have failed. Early help is far more effective than waiting until the consequences become severe.

\textbf{Contact your local emergency services for:}
\begin{itemize}
    \item Overdose: severe drowsiness, slowed or stopped breathing, or unresponsiveness
    \item Severe withdrawal, especially seizures, confusion, or delirium
    \item Suicidal thoughts or self-harm
    \item Any life-threatening reaction to a drug or alcohol
\end{itemize}

\section*{Diagnosis and Investigations}
\begin{itemize}
    \item \textbf{Clinical interview:} a detailed, non-judgemental history of substance use, its harms, and previous attempts to stop
    \item \textbf{Screening questionnaires:} such as the AUDIT for alcohol or the DAST for other drugs
    \item \textbf{Physical examination:} to look for signs of intoxication, withdrawal, or organ damage
    \item \textbf{Blood tests:} liver function, blood counts, and alcohol or drug levels
    \item \textbf{Mental health assessment:} to identify depression, anxiety, or other conditions
    \item \textbf{Infectious disease screening:} hepatitis B and C, and HIV, for people who inject drugs
\end{itemize}

\section*{Management}
Treatment is most effective when tailored to the person and their circumstances:

\begin{itemize}
    \item \textbf{GP care:} a first point of contact and coordinator of treatment
    \item \textbf{Counselling:} cognitive behaviour therapy and motivational interviewing
    \item \textbf{Medicines:} such as nicotine replacement therapy, buprenorphine or methadone for opioid dependence, and medicines that support abstinence from alcohol
    \item \textbf{Withdrawal support:} supervised detoxification, especially for alcohol and benzodiazepines
    \item \textbf{Rehabilitation and support groups:} residential programmes and mutual support groups such as Alcoholics Anonymous or Narcotics Anonymous
    \item \textbf{Family involvement:} support and education for partners and relatives
\end{itemize}

\section*{Complications}
\begin{itemize}
    \item Overdose and death
    \item Hepatitis B, hepatitis C, HIV, and other infections
    \item Liver, heart, and lung damage, and some cancers
    \item Mental illness, including depression and suicide
    \item Accidents and injuries while intoxicated
    \item Relationship breakdown, financial loss, and legal problems
\end{itemize}

\section*{Prognosis and Outcomes}
Addiction is a chronic condition, and recovery often involves several attempts. Many people do recover, particularly with treatment, support, and time. Relapse is common and should be seen as part of the process rather than as a failure.

With good support, people with addiction can regain control, rebuild relationships, and return to work or study. Outcomes are best when the person engages with a combination of medical treatment, counselling, and ongoing social support, and when any underlying mental health problems are also addressed.

\section*{Prevention and Self-Care}
\begin{itemize}
    \item Learn about the risks of alcohol and other drugs
    \item Set clear limits on drinking and avoid binge drinking
    \item Keep medicines secure and take them only as prescribed
    \item Build strong social connections and healthy ways to manage stress
    \item Seek help early for mental health problems
    \item If you are in recovery, avoid high-risk situations and people
    \item Talk openly with children and young people about drugs and alcohol
\end{itemize}

\section*{Sources and Further Reading}
The following organisations provide reliable information on addiction and its treatment.

\begin{itemize}
    \item NHS. \textit{Information on drug addiction and alcohol dependence.} https://www.nhs.uk/conditions/
    \item World Health Organization. \textit{Resources on substance use and dependence.} https://www.who.int/
    \item Centers for Disease Control and Prevention. \textit{Information on substance use and health.} https://www.cdc.gov/
    \item NICE. \textit{Clinical guidance on drug misuse and alcohol dependence.} https://www.nice.org.uk/
    \item Mayo Clinic. \textit{Guide to drug addiction and treatment options.} https://www.mayoclinic.org/
\end{itemize}
